\DocumentMetadata{
  pdfversion=2.0,
  pdfstandard=ua-2,
  lang=en-US,
  tagging=on,
  tagging-setup={math/setup=mathml-SE,tabsorder=structure}
}
\documentclass[11pt,letterpaper]{article}
\usepackage[margin=0.68in,includefoot]{geometry}
\usepackage{newcomputermodern}
\usepackage{amsmath,amscd,mathtools}
\usepackage{tikz,adjustbox}
\providecommand{\square}{\mdlgwhtsquare}
\usepackage[hidelinks]{hyperref}
\hypersetup{
  pdftitle={Numerical Analysis PhD Exam, May 2005},
  pdfauthor={University of Florida Department of Mathematics},
  pdfsubject={Graduate mathematics examination},
  pdfdisplaydoctitle=true,
  bookmarksopen=true
}
\setcounter{secnumdepth}{0}
\setlength{\parindent}{0pt}
\setlength{\parskip}{0.28em}
\setlength{\emergencystretch}{3em}
\setlength{\leftmargini}{2.15em}
\setlength{\leftmarginii}{2.65em}
\setlength{\leftmarginiii}{3em}
\pagestyle{empty}
\raggedbottom
\allowdisplaybreaks
\NewTaggingSocketPlug{block/list/label}{examproblem}{%
  \tagstructbegin{tag=Lbl}\tagstructend%
  \tagstructbegin{tag=\LBody}%
  \tagstructbegin{tag=\ProblemHeadingTag,title-o={\ProblemHeadingText}}%
  \tagmcbegin{tag=\ProblemHeadingTag,actualtext={\ProblemHeadingText}}%
  #2%
  \tagmcend%
  \tagstructend%
}
\newenvironment{examproblems}[2]{%
  \def\ProblemHeadingTag{#2}%
  \AssignTaggingSocketPlug{block/list/label}{examproblem}%
  \begin{enumerate}%
  \setcounter{enumi}{#1}%
  \renewcommand{\labelenumi}{\textbf{\theenumi.}}%
  \setlength{\topsep}{0.35em}%
  \setlength{\partopsep}{0pt}%
  \setlength{\itemsep}{0.48em}%
  \setlength{\parsep}{0.18em}%
}{\end{enumerate}}
\newenvironment{examsubparts}{%
  \AssignTaggingSocketPlug{block/list/label}{default}%
  \begin{enumerate}%
  \setlength{\topsep}{0.18em}%
  \setlength{\partopsep}{0pt}%
  \setlength{\itemsep}{0.16em}%
  \setlength{\parsep}{0.1em}%
}{\end{enumerate}}
\newenvironment{examinstructions}{%
  \par\begingroup\itshape
}{%
  \par\endgroup\vspace{0.18em}
}
\makeatletter
\renewcommand\subsection{\@startsection{subsection}{2}{0pt}%
  {-1.25ex plus -.25ex minus -.1ex}%
  {0.55ex plus .1ex}%
  {\normalfont\large\bfseries}}
\renewcommand{\@maketitle}{%
  \newpage\null\vskip -1em%
  {\footnotesize\bfseries
    \href{https://www.ufl.edu/}{University of Florida}, \href{https://math.ufl.edu/}{Department of Mathematics}\par}%
  \vskip -0.5em%
  \tagstructbegin{tag=H1,title={Numerical Analysis PhD Exam, May 2005}}%
  \tagpdfparaOff%
  \tagmcbegin{tag=H1}%
    {\LARGE\bfseries\href{https://gma.math.ufl.edu/exams/numerical-analysis-phd/}{Numerical Analysis PhD Exam}, May 2005\par}%
  \tagmcend%
  \tagpdfparaOn%
  \tagstructend%
  \vskip 0.25em%
  \tagmcbegin{artifact}%
    \hrule height 1pt%
  \tagmcend%
  \vskip 1em%
}
\makeatother
\title{Numerical Analysis PhD Exam, May 2005}
\author{}
\date{}
\begin{document}
\maketitle
\thispagestyle{empty}
\begin{examinstructions}
Do any eight of the ten problems below.
\end{examinstructions}
\begin{examproblems}{0}{H2}
\def\ProblemHeadingText{Problem 1.}
\item
This problem has three parts.
\begin{examsubparts}
\item[\textnormal{(a)}]
Consider the matrix \(A=uv^*\), where~\(u\) and \(v\in\mathbb{C}^n\). Under what condition on~\(u\) and~\(v\) is~\(A\) a projector?
\item[\textnormal{(b)}]
Show that the Householder matrix \(H=I-2ww^*\), where \(\lVert w\rVert=1\), is unitary.
\item[\textnormal{(c)}]
Given a vector \(x\in\mathbb{C}^n\) and an integer~\(k\) with \(1<k<n\), derive a formula for a Householder matrix with the property that \((Hx)_i=0\) for \(i>k\) and \((Hx)_i=x_i\) for \(i<k\). Be sure to choose the signs so that the formula is numerically stable.
\end{examsubparts}
\def\ProblemHeadingText{Problem 2.}
\item
This problem has three parts.
\begin{examsubparts}
\item[\textnormal{(a)}]
Given a matrix \(A\in\mathbb{C}^{m\times n}\) with~\(n\) linearly independent columns, derive the formula (the normal equation) for the least squares solution to the possibly overdetermined linear system \(Ax=b\).
\item[\textnormal{(b)}]
Given a matrix \(A\in\mathbb{C}^{m\times n}\) with~\(m\) linearly independent rows, derive the formula (the analogue of the normal equation) for the minimum norm solution \(Ax=b\).
\item[\textnormal{(c)}]
Suppose you are given the singular value decomposition (SVD) \(A=U\Sigma V^*\). Derive a single formula for the least squares/minimum norm solution to \(Ax=b\) which is expressed in terms of~\(U\), \(\Sigma\), and~\(V\).
\end{examsubparts}
\def\ProblemHeadingText{Problem 3.}
\item
Consider the invertible linear system \(Ax=b\). Derive a tight bound for the condition number with respect to perturbations in~\(b\). If \(A(x+\delta x)=b+\delta b\) is the perturbed system, then with the 2-norm, for which choice of~\(b\) and for which perturbation \(\delta b\) is this bound tight?
\def\ProblemHeadingText{Problem 4.}
\item
Let
\[
C=\begin{bmatrix}
-a_1&-a_2&\cdots&-a_{n-1}&-a_n\\
1&0&\cdots&0&0\\
0&1&\ddots&\vdots&\vdots\\
\vdots&\ddots&\ddots&0&\vdots\\
0&\cdots&0&1&0
\end{bmatrix},
\qquad
e=\begin{bmatrix}1\\0\\\vdots\\\vdots\\0\end{bmatrix}.
\]
\begin{examsubparts}
\item[\textnormal{(a)}]
What is the characteristic polynomial of~\(C\)?
\item[\textnormal{(b)}]
Consider any \(A\in\mathbb{C}^{n\times n}\) and \(b\in\mathbb{C}^n\). Let the characteristic polynomial of~\(A\) be
\[
p(\lambda)=\lambda^n+a_1\lambda^{n-1}+a_2\lambda^{n-2}+\cdots+a_{n-1}\lambda+a_n,
\]

\[
K=[b,Ab,A^2b,\ldots,A^{n-1}b],
\qquad
R=\begin{bmatrix}
1&a_1&a_2&\cdots&a_{n-1}\\
0&1&a_1&\cdots&a_{n-2}\\
\vdots&\ddots&\ddots&\ddots&\vdots\\
\vdots&&\ddots&\ddots&a_1\\
0&\cdots&\cdots&0&1
\end{bmatrix}.
\]
 Show that if~\(K\) is invertible, then
\[
(KR)^{-1}A(KR)=C
\qquad\text{and}\qquad
(KR)^{-1}b=e.
\]
 (Suggestion: Prove that \(A(KR)=(KR)C\) using \(p(A)=0\).)
\end{examsubparts}
\def\ProblemHeadingText{Problem 5.}
\item
Let \(A=QR\) be the factorization of~\(A\) into the product of a unitary matrix and a triangular matrix. Suppose that the columns of~\(A\) are linearly independent. Show that \(|r_{kk}|\) is the distance from the~\(k\)-th column of~\(A\) to the linear space spanned by the first \(k-1\) columns of~\(A\).
\def\ProblemHeadingText{Problem 6.}
\item
Show that \(x=\pi+0.5\sin(x/2)\) has a unique solution~\(\alpha\) in~\(\mathbb{R}\). Give a fixed point iteration by which one can compute an approximation to~\(\alpha\) numerically and prove that it converges at order one.
\def\ProblemHeadingText{Problem 7.}
\item
The Lagrange interpolation problem is to find a polynomial \(p_n(x)\) of degree at most~\(n\) such that
\[
p_n(x_i)=y_i,\qquad i=0,\ldots,n,
\]
 where \(x_0,\ldots,x_n\) are distinct real points and \(y_0,\ldots,y_n\) are given data. Prove that this problem always has a unique solution.
\def\ProblemHeadingText{Problem 8.}
\item
Consider the quadrature that approximates
\[
\int_a^b f(x)\,dx
\]
 by
\[
I_3(f):=\int_a^b P_3(x)\,dx,
\]
 where \(P_3(x)\) is the cubic Hermite interpolant to \(f(x)\) at the points \(x=a,b\). Obtain an error formula for this approximation. (This is called the corrected trapezoid rule.)
\def\ProblemHeadingText{Problem 9.}
\item
Consider multistep methods of the form
\[
y_{n+1}=y_{n-q}+h\sum_{j=-1}^{p}b_jf(x_{n-j},y_{n-j})
\]
 with \(q\geq1\). Show that such methods do not satisfy the strong root condition. Find an example with \(q=1\) that is relatively stable.
\def\ProblemHeadingText{Problem 10.}
\item
Let \(S_n^1\) denote the space of linear splines based on knots \(a=t_0<t_1<\cdots<t_n=b\), i.e.,
\[
S_n^1=\left\{v:v|_{[t_j,t_{j+1}]}\text{ is linear for all }j=0,1,\ldots,n-1\text{ and }v\text{ is continuous on }[a,b]\right\}.
\]
 Let \(s_f\in S_n^1\) interpolate a continuous function~\(f\) at the knots.
\begin{examsubparts}
\item[\textnormal{(a)}]
Prove that \(\lVert s_f\rVert_\infty\leq\lVert f\rVert_\infty\). (Here \(\lVert v\rVert_\infty=\max_{x\in[a,b]}|v(x)|\).)
\item[\textnormal{(b)}]
Prove that \(\lVert f-s_f\rVert_\infty\leq2\lVert f-s\rVert_\infty\) for any \(s\in S_n^1\).
\item[\textnormal{(c)}]
Consider the error in best approximation to~\(f\) from \(S_n^1\) and the error in approximating~\(f\) by its linear interpolant \(s_f\). If one goes to zero, does the other go to zero?
\end{examsubparts}
\end{examproblems}
\end{document}
